\documentclass[11pt]{article}
\usepackage{fontspec,amsmath,amssymb,amsthm,hyperref}
\usepackage[margin=1in]{geometry}
\setlength{\emergencystretch}{3em}
\title{Jointly measurable density and relative score}
\author{LangevinLean mathematical exposition}
\date{}
\begin{document}\maketitle
\noindent Source module: \texttt{UnadjustedLangevin/LSI/GaussianMeasurability.lean}. Compile this proof document with XeLaTeX or LuaLaTeX. The accompanying HTML pairs the mathematical argument with the exact Lean statements and records the scope of the formalization.


Spacetime integration requires joint measurability in time and state, rather than measurability separately at each fixed time. This module constructs a single measurable extension of the actual positive-time relative score and proves that it agrees with the gradient used in Fisher information.


\section{Integrate a jointly measurable Gaussian kernel}


The density is a nonnegative parameter integral \(p_t(y)=\int K_t(x,y)d\mu(x)\). The normalizer, affine center, norm, and exponential are measurable in \((t,y,x)\). The parameterized nonnegative integral is therefore measurable. Passing from the extended integral to its real value preserves measurability. Multiplication by \(Z e^{U(y)}\) gives the joint measurability of \(r_t(y)\).



The definitions cover all real times using Lean’s total real operations; only positive times are used for density and score identities. This avoids defining a different measurable representative for each interval.



\noindent Lean declarations: \texttt{interpolationLebesgueDensity\_joint\_measurable}, \texttt{interpolationRelativeDensity\_joint\_measurable}.

\section{Measurability of the spatial derivative}


The derivative kernel is a continuous linear functional proportional to the inner product with the Gaussian residual:

\[D_yK_t(x,y)[v]=-\frac{K_t(x,y)}{2t}\langle y-x+t\nabla U(x),v\rangle.\]

The inner-product-to-dual map is continuous, so this is jointly strongly measurable as a continuous-linear-map-valued function. Its Bochner integral in \(x\), denoted \(A_t(y)\), remains strongly measurable. At positive times the earlier differentiation theorem identifies \(A_t=D p_t\).



\noindent Lean declarations: \texttt{interpolationLebesgueDensityDeriv\_joint\_stronglyMeasurable}.

\section{One relative-score field for all spacetime integrals}


Define

\[\mathcal S(t,y)=\nabla U(y)+p_t(y)^{-1}\mathcal R^{-1}(A_t(y)),\]

where \(\mathcal R^{-1}\) is the inverse Riesz identification of the continuous dual with the Hilbert space. Continuous scalar multiplication, inversion as a measurable real function, and the previous measurability results make \(\mathcal S\) jointly strongly measurable.



For \(t>0\), \(p_t>0\), so the log chain rule and \(r_t=Z e^U p_t\) give

\[\mathcal S(t,y)=\nabla U(y)+\nabla\log p_t(y)=\nabla\log r_t(y).\]

Thus the measurable field is the actual relative score wherever it is used. Its auxiliary extension at nonpositive time introduces no analytic premise into the proof.



\noindent Lean declarations: \texttt{interpolationRelativeScoreField}, \texttt{interpolationRelativeScoreField\_stronglyMeasurable}, \texttt{interpolationRelativeScoreField\_eq}.

\end{document}
